
\chapter[APÊNDICE A -- PLANO DE TRABALHO: DIVISÃO DE RESPONSABILIDADES]{\fbox{\parbox{0.9\textwidth}{\centering APÊNDICE A \textendash{} PLANO DE TRABALHO: DIVISÃO DE RESPONSABILIDADES}}} \label{ap:plano}

Este apêndice detalha a divisão de responsabilidades entre os dois autores, conforme a pré-proposta aprovada e o plano de trabalho aprovado pelo Colegiado do Curso de Engenharia de Computação do Centro Federal de Educação Tecnológica de Minas Gerais, Unidade Leopoldina.

\section{Atividades Desenvolvidas em Regime Colaborativo}

As atividades listadas a seguir, de natureza estruturante ou situadas na fronteira entre os dois subsistemas, foram conduzidas de forma colaborativa por Victor de Souza Vilela da Silva e Igor Lamoia Queiroz:

\begin{itemize}
    \item Revisão bibliográfica sobre interpretadores, compiladores e ferramentas educacionais correlatas;
    \item Definição dos requisitos funcionais e não funcionais e da arquitetura geral do sistema;
    \item Implementação e validação do núcleo do interpretador, abrangendo as fases de análise léxica, análise sintática, geração de código intermediário e execução, conduzidas conjuntamente pelos dois autores;
    \item Definição dos contratos de comunicação entre o núcleo do interpretador e a plataforma web e a integração entre os módulos;
    \item Revisão mútua de código e tomada de decisões arquiteturais nas etapas de refatoração e evolução do sistema.
\end{itemize}

\section{Eixos de Aprofundamento Individual}

A partir da base técnica construída em conjunto, cada autor aprofundou sua análise acadêmica em eixos específicos, conforme previsto na pré-proposta.

\subsection{Victor de Souza Vilela da Silva (esta monografia)}

\begin{itemize}
    \item Detalhamento das fases de análise léxica e do sistema de varredura por escâneres especializados (\textit{factory pattern});
    \item Descrição aprofundada da análise sintática por descida recursiva e da correspondência entre as produções do analisador e a gramática formal da linguagem Java-{}-;
    \item Detalhamento e documentação da geração de código de três endereços: instrução a instrução, tabela de temporários e tabela de rótulos;
    \item Análise aprofundada dos mecanismos de injeção dinâmica de regras gramaticais no interpretador, que permitem ao núcleo reconfigurar o vocabulário e as regras estruturais a cada execução;
    \item Análise do funcionamento do interpretador, da tabela de símbolos e do despacho das chamadas de sistema de entrada e saída;
    \item Estudos de caso centrados na validade técnica do núcleo de tradução e no comportamento do interpretador frente a programas de teste elaborados para cobrir as construções da linguagem.
\end{itemize}

\subsection{Igor Lamoia Queiroz (documento complementar)}

\begin{itemize}
    \item Detalhamento da arquitetura da aplicação web (Next.js, Pages Router) e da integração com o Monaco Editor para edição de código com realce de sintaxe dinâmico;
    \item Documentação do assistente de personalização da linguagem (\textit{wizard}), incluindo o modelo de etapas, a lógica de validação preventiva dos delimitadores de bloco e o mecanismo de pré-visualização em tempo real;
    \item Análise da experiência de uso da plataforma, dos contextos React que governam o estado da aplicação e das mediações pedagógicas oferecidas pela interface;
    \item Descrição da infraestrutura de submissão e correção automatizada de exercícios e do sistema de gerenciamento de turmas;
    \item Estudos de caso centrados na usabilidade, no engajamento dos estudantes e na retroalimentação coletada em ambiente educacional.
\end{itemize}

\chapter[APÊNDICE B -- DOCUMENTAÇÃO DA API]{\fbox{\parbox{0.9\textwidth}{\centering APÊNDICE B \textendash{} DOCUMENTAÇÃO DA API}}} \label{ap:api}

Este apêndice reúne a relação dos recursos expostos pela API da plataforma. A documentação completa e interativa é gerada automaticamente pelo \textit{framework} FastAPI e está disponível \textit{online} nas interfaces Swagger UI (\url{https://api.dynamicinterpreter.com.br/docs}) e ReDoc (\url{https://api.dynamicinterpreter.com.br/redoc}). A Tabela~\ref{tab:api-recursos} sintetiza os grupos de recursos e suas responsabilidades.

\begin{table}[H]
    \centering
    \caption{Grupos de recursos da API da plataforma.}
    \label{tab:api-recursos}
    \begin{tabular}{|>{\raggedright\arraybackslash}p{4.3cm}|p{8.3cm}|}
        \hline
        \textbf{Recurso} & \textbf{Responsabilidade} \\
        \hline
        \texttt{POST /api/}\allowbreak\texttt{submissions/}\allowbreak\texttt{validate} & Executa o ciclo completo de compilação e interpretação do código submetido e o avalia contra os casos de teste do exercício (rota do pacote \texttt{ide}). \\
        \hline
        \texttt{/auth} & Registro e \textit{login} de usuários, com emissão de \textit{tokens} JWT. \\
        \hline
        \texttt{/organizations} & Gestão das instituições de ensino cadastradas. \\
        \hline
        \texttt{/users} & Consulta e atualização de perfil, com controle de acesso por \texttt{role}. \\
        \hline
        \texttt{/classes} & Criação, edição, encerramento e arquivamento de turmas; matrícula via \texttt{/classes/join}. \\
        \hline
        \texttt{/exercises} & Criação e manutenção de enunciados e respectivos casos de teste. \\
        \hline
        \texttt{/exercise-lists} & Composição, ordenação e atribuição de pesos aos exercícios de uma lista. \\
        \hline
        \texttt{/class-}\allowbreak\texttt{exercise-}\allowbreak\texttt{lists} & Publicação de listas para turmas, com prazo, nota total e mínimo de exercícios. \\
        \hline
        \texttt{/submissions} & Armazenamento e recuperação das submissões dos estudantes, com pontuação e retroalimentação. \\
        \hline
    \end{tabular}
    \legend{Fonte: elaborado pelo autor.}
\end{table}

\chapter[APÊNDICE C -- CASOS DE TESTE AUTOMATIZADOS]{\fbox{\parbox{0.9\textwidth}{\centering APÊNDICE C \textendash{} CASOS DE TESTE AUTOMATIZADOS}}} \label{ap:testes}

Este apêndice resume a suíte de testes automatizados do núcleo do compilador, construída com o \textit{framework} Vitest. A suíte reúne \textbf{215 casos} distribuídos em 22 arquivos, organizados de modo a refletir as camadas internas do projeto, conforme a Tabela~\ref{tab:testes-resumo}. O conjunto completo dos casos e o respectivo relatório de execução estão disponíveis no repositório do projeto (\url{https://github.com/igorlamoia/ts-compilator-for-java}, em \texttt{packages/compiler/src/tests}); à data desta qualificação, a totalidade dos casos é aprovada.

\begin{table}[H]
    \centering
    \caption{Distribuição dos casos de teste por camada do compilador.}
    \label{tab:testes-resumo}
    \begin{tabular}{|>{\raggedright\arraybackslash}p{3.4cm}|c|c|p{6.2cm}|}
        \hline
        \textbf{Camada} & \textbf{Arquivos} & \textbf{Casos} & \textbf{Escopo} \\
        \hline
        Tokens & 1 & 8 & Estabilidade dos identificadores numéricos e dos agrupamentos por categoria. \\
        \hline
        Análise léxica & 9 & 56 & \textit{Scanners} em isolamento: comentários, literais numéricos e de cadeia, terminador configurável, delimitadores de bloco, operadores em forma de palavra e geração de \textit{tokens} de indentação. \\
        \hline
        Análise sintática e semântica & 9 & 138 & Aceitação ou rejeição de programas sob distintos modos de tipagem e de bloco, variantes do \texttt{switch} e verificação de tipos. \\
        \hline
        Interpretação & 2 & 13 & Execução do código de três endereços e sessões de depuração. \\
        \hline
        \textbf{Total} & \textbf{22} & \textbf{215} & \\
        \hline
    \end{tabular}
    \legend{Fonte: elaborado pelo autor.}
\end{table}
